% 2010 年辽宁卷理科第 22 题：△ABC 的外接圆，角平分线 AD 的延长线交圆于 E。
% 来源 PDF：普通高考/2010/2010辽宁理.pdf，SHA-256 ba46f70f9d718652c24d14b28a0b473b768e864248ff0a9f1feb2374aec708f9。
% 原页：第 2 页（A3 三栏）第 2 栏下部，书页码 612。
% 重绘依据：照原页——BC 为水平弦（本图中恰为直径），A 在圆上偏左上方，E 在圆上偏下方；
%   AD 与 AE 是同一条直线，D 为它与 BC 的交点；线段 AB、AC、BC、AE、BE 均为实线。
\begin{tikzpicture}[scale=1.15, semithick, line join=round]
  \coordinate (O) at (0,0);
  \coordinate (B) at (180:2);
  \coordinate (C) at (0:2);
  \coordinate (A) at (123:2);
  \coordinate (E) at (266:2);
  \coordinate (D) at (intersection of A--E and B--C);

  \draw (O) circle[radius=2];

  \draw (A) -- (B);
  \draw (A) -- (C);
  \draw (B) -- (C);
  \draw (A) -- (E);
  \draw (B) -- (E);

  \node[above, inner sep=2pt] at (A) {$A$};
  \node[left, inner sep=2pt] at (B) {$B$};
  \node[right, inner sep=2pt] at (C) {$C$};
  \node[below right=1pt and 1pt, inner sep=1pt] at (D) {$D$};
  \node[below, inner sep=2pt] at (E) {$E$};
\end{tikzpicture}
